\documentclass[12pt]{article}
\usepackage[a4paper, margin=1in]{geometry}
\usepackage{amsmath, minted, enumitem, cmbright, hyperref}
\renewcommand{\thesection}{\Alph{section}}
\renewcommand{\thesubsection}{\thesection.\arabic{subsection}}
\renewcommand{\t}{\texttt}
\begin{document}
\section{Short Questions}
\textit{Note that the original problem is written in Java but I include the answers for the sake of completeness.}
\subsection{Java Code Comprehension and Analysis}
\begin{enumerate}
\item \t{3}
\item \t{0}
\item The number of contiguous subarrays of length $m$ that contain at least $2$ even numbers.
\item There are $n-m+1$ windows each of which scans $m$ elements. Thus the tightest time complexity is $O(mn)$.
\item Use a sliding window, which is $O(n)$. Code omitted.
\end{enumerate}
\subsection{Big-O Time Complexity Analysis}
\begin{enumerate}
\item Can be false. Big-O does not give exact operation counts.
\item Can be false. For instance $O(N\log N)$ is already optimal for (comparison-based) sorting and cannot be improved to $O(N)$.
\item Always true. For sufficiently large $N$ we have $N\log N<N^2$, so $O(N\log N)\subset O(N^2)$.
\item Can be false. It depends on the operations used.
\item Always true. The algorithm is $O(N\log N)$ so it takes at most $cN\log N$ operations for some constant $c$, but for any $c$ and sufficiently large $N$ we have $cN\log N<N^2$.
\end{enumerate}
\subsection{Another Merge Sort Variant}
\begin{enumerate}
\item Slower. We are changing an $O(n)$ merge to an $O(n\log n)$ sort, so it must be (asymptotically) slower.
\item $O\left(n\log^2 n\right)$. We can show this via \href{https://en.wikipedia.org/wiki/Master_theorem_(analysis_of_algorithms)#Generic_form}{master's theorem (case 2)}. Indeed we have $T(n)=2T(n/2)+O(n\log n)$ so $T(n)\in O\left(n\log^2 n\right)$.
\end{enumerate}
\subsection{Analyze These Statements}
\begin{enumerate}
\item Incorrect. We have $\sum_{i=1}^n\lfloor n/i\rfloor\in O(n\log n)$.
\item Incorrect. It has worst time complexity $O\left(n^2\right)$ which is beaten by e.g., merge sort with $O(n\log n)$.
\item Correct. $O(n\log n)$ assumes that comparisons are $O(1)$, but comparisons between strings of length $m$ can be $O(m)$ in the worst case. Thus sorting can be $O(mn\log n)$ which is more than $O(n^2)$ if e.g., $m=n^2$.
\item Correct. We can achieve $O(n\log n)$ via a modified version of merge sort.
\item Incorrect. While both are (possibly amortised) $O(1)$, \t{ArrayList} can be faster due to better cache locality.
\item Incorrect. As the maximum number of elements is already known beforehand, we can avoid resizing and use a circular array, which is $O(1)$. Again \t{ArrayList} can be faster due to better cache locality.
\item Incorrect. For instance, if all the integers are equal.
\item Incorrect. Since elements are from only 26 possible letters, we can use a hash array and solve it in $O(n+m)$ with something like
\begin{minted}{java}
boolean[] s = new boolean[26];
for (Character a : La)
    s[a - 'A'] = true;
for (Character b : Lb)
    if (s[b - 'A']) return true;
return false;
\end{minted}
\item Incorrect. It is always in a leaf but not necessarily at the rightmost.
\item Correct. We can just sort the queue once after all the enqueue operations are performed, which is also $O(n\log n)$.
\end{enumerate}
\subsection{An Algorithm that uses Priority Queue}
\textit{Note that this is just problem C from CS2040S 2023/24 S1 midterm.}
\begin{enumerate}
\item \t{1+3=4
4+5=9
6+9=15
10+15=25}
\item \t{12+16=28}
\item $\sum_{i=2}^nO(i\log i)=O\left(n^2\log n\right)$.
\item Use a min-heap which is $O(n\log n)$. Code omitted. (This is actually used in \href{https://en.wikipedia.org/wiki/Huffman_coding}{Huffman coding} which is covered in CS3230.)
\end{enumerate}
\end{document}